%% kapa: exam=1 round=33 year=2022 period=2 form=A subject=appraisal_law subject_ko=감정평가관계법규 seq=23 number=23
%% source: Q-Net 감정평가사 기출문제 (cst00309) artlSeq=5209997 file=raw/1008/5209997/2022년도 제33회 감정평가사 1차시험 2교시(감정평가 관계법규, 회계학).pdf
%% conversion: kordoc 4.24.51 -> tools/gen1.py (auto); review=proofread-sonnet; flags=none
\begin{question}{23}
국유재산법령상 지식재산에 관한 설명으로 옳지 않은 것은?
\begin{choices}
  \item 「식물신품종 보호법」에 따른 품종보호권은 지식재산에 해당한다.
  \item 지식재산을 대부 받은 자는 해당 중앙관서의 장등의 승인을 받아 그 지식재산을 다른 사람에게 사용·수익하게 할 수 있다.
  \item 상표권의 사용료를 면제하는 경우 그 면제기간은 5년 이내로 한다.
  \item 저작권등의 사용허가를 받은 자는 해당 지식재산을 관리하는 중앙관서의 장등의 승인을 받아 그 저작물의 개작을 할 수 있다.
  \item 지식재산의 사용허가등의 기간을 연장하는 경우 최초의 사용허가등의 기간과 연장된 사용허가등의 기간을 합산한 기간은 5년을 초과하지 못한다.
\end{choices}
\end{question}
\answer{3}
